\section{The exact weight-nine arithmetic normalization}
\label{sec:weight-nine-bounds}

The next odd coefficient gives another exact arithmetic quotient.
Use $G=D_8H$, $J_7=D_7\cap H$ and $J_8=D_8\cap H$.
For $g=(1,f)\in D_7$, put
\[
 m_9(g)=[S^8T]f.
\]
We prove that its image on each of $D_7,D_8,J_7,J_8$ is
$512\mathbf Z_2$. Killing this coefficient on $D_8$ therefore
preserves the relative abelian obstruction and leaves only words
starting in degree ten. The uniform odd theorem supplies the integral bound. We correct an
arithmetic value by the degree-eight commutator and then determine the
whole leading line.

\subsection{An additive coefficient and an arithmetic value}

For $g\in M=\ker\lambda$, write $a_j=[S^jT]f$. The logarithmic
Alexander beta coefficient
\[
 c_9(g)=[u^8v]\log B_g(e^u,e^v)
\]
is a continuous homomorphism on $M$ and has conjugation weight
$\lambda_2^9$. The boundary formula for the beta cocycle gives
\[
 c_9=[u^8]\sum_{j\ge1}a_j(e^u-1)^j.
\]
Using $a_1=0$, $a_3=-a_2$, and the finite exponential expansion gives
\begin{equation}\label{eq:weight-nine-beta-formula}
 \begin{split}
 \tau_9:=20160c_9={}&20160a_8+70560a_7+95760a_6\\
                  &+63000a_5+20412a_4-2771a_2.
 \end{split}
\end{equation}
This is an integral homomorphism on all of $M$. Explicitly, the
coefficients of $a_2,a_4,a_5,a_6,a_7,a_8$ in $c_9$ are
\[
 -\frac{2771}{20160},\quad\frac{81}{80},\quad
 \frac{25}{8},\quad\frac{19}{4},\quad\frac72,\quad1.
\]
These numbers can also be checked from
$[u^8](e^u-1)^j=j!\genfrac{\{}{\}}{0pt}{}{8}{j}/8!$.
On $D_7$ the lower depth-one terms vanish: the possible degree-eight
bracket has no depth-one component. Consequently $c_9=m_9$ already
on $D_7$, and in particular on $D_8$.

\begin{proposition}\label{prop:weight-nine-arithmetic-value}
There is $w_9\in J_8$ such that
\[
 m_9(w_9)=512\cdot12746612385\,\tau_9(\sigma),
 \qquad \tau_9(\sigma)\in\mathbf Z_2^\times,
\]
where $\sigma$ is the arithmetic element used in the preceding
normalizations. Thus $512\mathbf Z_2\subseteq m_9(J_8)$.
\end{proposition}

\begin{proof}
Choose $\sigma\in H\cap M$ with $\rho_2(\sigma)=2$ and all the
odd-prime Kummer coordinates zero, and choose $h\in H$ with
$\lambda_2(h)=3$. Put $t=\lambda(h)$ and
\[
 \Delta_j(z)=hzh^{-1}z^{-t^j},\qquad
 z_1=\Delta_1(\sigma),\quad z_3=\Delta_3(z_1),\quad
 z_5=\Delta_5(z_3),\quad z_7=\Delta_7(z_5).
\]
All powers are the continuous profinite powers. The previously
proved weight-one, three, five and seven conjugation laws show
successively that $z_1\in D\cap H$, $z_3\in J_3$, $z_5\in J_5$,
and $z_7\in J_7$. By the weight-nine law,
\[
 c_9(z_7)=\frac{\prod_{j\in\{1,3,5,7\}}(3^9-3^j)}{20160}
                          \tau_9(\sigma)
        =512\cdot12746612385\,\tau_9(\sigma).
\]
The integer multiplying $512$ is odd; the full numerator divided
by $20160$ is $2^9 3^{14}\cdot5\cdot13\cdot41$.
Also $m_3(\sigma)$ is odd, and all other coefficients in
\eqref{eq:weight-nine-beta-formula} are even. Thus
$\tau_9(\sigma)$ is odd.

Let $w_8\in J_7$ be the primitive arithmetic commutator of the
preceding section. Its $m_8$ value is a unit. Choose
$u\in\widehat{\mathbf Z}$ whose $2$-adic component is
$m_8(z_7)/m_8(w_8)$, and set
\[
 w_9=z_7w_8^{-u}.
\]
Then $w_9\in J_8$. Since $w_8$ is a commutator of two elements of
$M$ and $c_9$ is a homomorphism on $M$, one has $c_9(w_8)=0$.
Therefore $m_9(w_9)=c_9(w_9)=c_9(z_7)$, as claimed. Compactness
of the arithmetic image supplies the entire ideal
$512\mathbf Z_2$.
\end{proof}

\subsection{The full leading line}

Use the standard Lyndon Lie and group commutators, with $S<T$ and
the longest proper Lyndon suffix. The degree-nine leading Lie
module of multiplier-one GT words that start in degree nine is contained
in $\mathbf Z_2\mathcal L_9$, where the primitive vector
$\mathcal L_9$ has the following coordinates:

{\small
\[
\begin{array}{cr@{\quad}cr@{\quad}cr}
w&[L_w]\mathcal L_9&w&[L_w]\mathcal L_9&w&[L_w]\mathcal L_9\\ \hline
SSSSSSSST&144&SSSTSTSST&-1118&SSTSTTTTT&-3686\\
SSSSSSSTT&-576&SSSTSTSTT&-5303&SSTTSSTTT&-1013\\
SSSSSSTST&1512&SSSTSTTST&-2951&SSTTSTSTT&2059\\
SSSSSSTTT&1344&SSSTSTTTT&4353&SSTTSTTST&5340\\
SSSSSTSST&-1440&SSSTTSSTT&-682&SSTTSTTTT&-3340\\
SSSSSTSTT&-346&SSSTTSTST&-11276&SSTTTSTST&4388\\
SSSSSTTST&-6322&SSSTTSTTT&479&SSTTTSTTT&1868\\
SSSSSTTTT&-2016&SSSTTTSST&-1762&SSTTTTSTT&3558\\
SSSSTSSST&360&SSSTTTSTT&-8109&SSTTTTTST&2434\\
SSSSTSSTT&-1598&SSSTTTTST&-6663&SSTTTTTTT&576\\
SSSSTSTST&3214&SSSTTTTTT&-1344&STSTSTSTT&3299\\
SSSSTSTTT&-2337&SSTSSTSTT&-89&STSTSTTTT&-5446\\
SSSSTTSST&2722&SSTSSTTST&-5496&STSTTSTTT&-1768\\
SSSSTTSTT&4724&SSTSSTTTT&-1949&STSTTTSTT&6452\\
SSSSTTTST&8823&SSTSTSSTT&291&STSTTTTTT&1368\\
SSSSTTTTT&2016&SSTSTSTST&1041&STTSTTTTT&2088\\
SSSTSSSTT&600&SSTSTSTTT&7472&STTTSTTTT&936\\
SSSTSSTST&2148&SSTSTTSTT&-4352&STTTTTTTT&-144\\
SSSTSSTTT&3259&SSTSTTTST&-2416&&\\
\end{array}
\]
}

Its first coordinate is $144$, and its nineteenth coordinate is
$3259$, an odd integer. Hence, if the leading vector is
$A\mathcal L_9$, then $A\in\mathbf Z_2$ and $m_9=144A$.

Here is complete finite data for the necessary leading line.
Collect group commutators of weights nine and ten, giving $155$
columns. Form $I_v,II_v$ as in the preceding section, now through
degree ten. Encode them by $2^{|v|}+\operatorname{bin}(v)$ and
$2048+2^{|v|}+\operatorname{bin}(v)$, respectively.
For the pentagon rows, let $d_{ij}$ be the derivation on the four
free letters $X_1,\ldots,X_4$ defined by
\[
 d_{ij}(X_i)=[X_i,X_j],\qquad
 d_{ij}(X_j)=[X_j,X_i],\qquad d_{ij}(X_k)=0\ (k\ne i,j).
\]
Use the five pairs and signs
\[
 (d_{12},d_{23}+d_{24},+),\quad
 (d_{13}+d_{23},d_{34},+),\quad
 (d_{23},d_{34},-),\quad
 (d_{12}+d_{13},d_{24}+d_{34},-),\quad
 (d_{12},d_{23},-).
\]
For a homogeneous Lie polynomial $P$, let $\mathcal P(P)$ be the
signed sum of its five evaluations on these pairs. This is the
leading Artin image of the pentagon, up to the common sign $(-1)^9$;
therefore it vanishes for every GT word under consideration.
Let $P_j$ be the coefficient on $X_1$ at the following five words,
respectively:
\[
 1111112324,\quad1111123324,\quad1111124324,\quad
 1111232324,\quad1111232424,
\]
where a digit $i$ denotes $X_i$. Give $P_j$ the row code $4096+j$.
Their entries on the degree-ten Hall columns are zero. The
coefficients on the degree-nine columns are determined by
$L_w(\mathcal U,\mathcal V)(X_1)$, with Lie bracket interpreted as
the commutator of derivations.

Let $A_R$ consist of the following $121$ rows, in the written order:
\[
\begin{aligned}
R=(&513,515,517,519,521,523,525,527,529,531,533,535,537,539,\\
&541,547,549,551,553,555,557,563,565,569,587,589,595,597,\\
&1025,1027,1029,1031,1033,1035,1037,1039,1041,1043,1045,1047,1049,1051,\\
&1053,1059,1061,1063,1065,1067,1069,1073,1075,1077,1079,1081,1083,1085,\\
&1093,1095,1097,1099,1101,1107,1109,1113,1115,1117,1125,1129,1131,1133,\\
&1139,1171,1173,1197,1203,1237,2563,2565,2569,2577,2583,2587,2603,2605,\\
&2613,3073,3075,3077,3079,3081,3083,3085,3089,3091,3093,3095,3097,3099,\\
&3101,3107,3109,3113,3115,3117,3119,3121,3123,3141,3145,3147,3149,3157,\\
&3159,3161,3163,3221,3243,3253,4097,4098,4101).
\end{aligned}
\]
Let $C$ be its first $121$ columns, and write
$q_j=\ell_j/144$, where $\ell_j$ is the $j$th coordinate of
$\mathcal L_9$ in lexicographic order. Define $T$ by its $55$ rows
$e_j-q_je_1$ for $2\le j\le56$. The exact finite identities are
\[
 \det A_{R,C}=22883210311680,\qquad
 T_C A_{R,C}^{-1}A_R=T.
\]
All entries are specified by the finite group-commutator and
geometric-series recursions through degree ten, and the displayed
derivations. Thus the identities give $h_j=q_jh_1$ for every
leading vector of a GT word. In particular it is a rational multiple
of $\mathcal L_9$, and its nineteenth coordinate makes the multiple
integral. This proves the asserted upper lattice without claiming
that arbitrary solutions of the necessary equations lift.

\subsection{The exact weight-nine normalization}
\label{subsec:weight-nine-exact}

The integral boundary measure forces the ninth coefficient into
its arithmetic lattice.

\begin{theorem}\label{thm:weight-nine-exact-image}
One has
\[
 m_9(D_8)=m_9(J_8)=512\mathbf Z_2.
\]
The same equality holds with $D_8,J_8$ replaced by $D_7,J_7$.
For $D_9=D_8\cap\ker m_9$ and $J_9=D_9\cap H$,
\[
 D_8=D_9J_8,\qquad G=D_9H,\qquad
 B\simeq D_9/\overline{[G,D_9]J_9}.
\]
The subgroup $D_9$ is proper, closed and normal in $G$, and every
word belonging to it starts in degree at least ten. Moreover,
\[
 (m_8,m_9):D_7\longrightarrow\mathbf Z_2\times512\mathbf Z_2
\]
is surjective already on $J_7$ and has kernel $D_9$.
\end{theorem}

\begin{proof}
The arithmetic word in Proposition~\ref{prop:weight-nine-arithmetic-value}
gives a value of exact valuation nine on $J_8$.
For the upper bound, use Lemma~\ref{lem:uniform-odd-coordinate-bound}
with $r=4$. On $D_7$ the coefficients $a_2,a_4,a_6$ vanish, so
$m_9=a_8\in512\mathbf Z_2$. Concretely, the polynomial
\[
 (t^2-1)(t^2-9)(t^2-25)(t^2-49)
 =t^8-84t^6+1974t^4-12916t^2+11025
\]
is divisible by $2^{15}$ on odd $t$. Its lower moments vanish,
and its eighth moment is $(1-2^8)\tau_9$, with
$\tau_9=20160m_9$ and $v_2(20160)=6$. This is exactly the required
valuation-nine bound. Compactness and the arithmetic witness
therefore prove equality of the two images.

The coefficient $m_9=c_9$ is already an additive, weight-nine
homomorphism on $D_7$. Given any element of $D_7$, multiply it by
a suitable profinite power of the primitive arithmetic
weight-eight commutator $w_8$ to make $m_8$ zero. This does not
change $c_9$, because every logarithmic beta homomorphism vanishes
on $w_8$. Consequently the same exact image equality holds on
$D_7$ and $J_7$. The two arithmetic words $w_8,w_9$ have respective
coordinate values
\[
 (m_8(w_8),0),\qquad(0,m_9(w_9)),
\]
where the first is a unit and the second has exact valuation
nine. They prove the asserted joint surjectivity.

The scalar conjugation laws imply normality of $D_9$. The
arithmetic value proves properness, and equality of the two images
gives $D_8=D_9J_8$, hence $G=D_9H$. The same compact gluing and
centralization argument as in the preceding normalizations gives
the displayed description of $B$. The leading-line calculation
shows that $m_9=0$ kills the whole degree-nine term, so the words
of $D_9$ start in degree at least ten.
\end{proof}

\subsection{The ninth character and its extension obstruction}

\begin{proposition}\label{prop:eta-nine-nonextension}
The character
\[
 \eta_9:D_7\longrightarrow\mathbf F_2,\qquad
 \eta_9(g)=m_9(g)/512\pmod2
\]
is continuous and $G$-invariant. Its restriction to $D_8$ has
nonzero arithmetic restriction and has no $G$-invariant extension
to $D_5$.
\end{proposition}

\begin{proof}
Integrality on $D_7$ and invariance follow from the exact image
and the weight-nine law. Recall the arithmetic construction
\[
 z_7=h z_5h^{-1}z_5^{-t^7}\in J_7,
 \qquad z_5\in J_5,
 \qquad w_9=z_7w_8^{-u}\in J_8,
\]
where $u_2=m_8(z_7)/m_8(w_8)$ and $m_8(w_8)$ is a unit.
The invariant character $\alpha_8$ on $D_5$ vanishes on $z_7$,
since $t$ is odd. Its restriction to $D_7$ is $m_8\bmod2$, so
$u_2$ is even.

Let $\chi:D_5\to\mathbf F_2$ be $G$-invariant. The same odd-power
argument gives $\chi(z_7)=0$. Also
$\chi(w_8^{-u})=-u_2\chi(w_8)=0$. Hence $\chi(w_9)=0$, whereas
$\eta_9(w_9)=1$ by its exact arithmetic valuation. This rules out
an invariant extension from $D_8$ to $D_5$.
\end{proof}

\begin{corollary}\label{cor:weight-nine-character-kernels}
The invariant quadratic-character kernels on restriction to $D_9$
are as follows:
\[
\begin{array}{c|l}
 \text{domain}&\text{kernel}\\ \hline
 D_8&\mathbf F_2\eta_9\\
 D_7&\mathbf F_2\eta_8\oplus\mathbf F_2\eta_9\\
 D_5&\mathbf F_2\eta_7\oplus\mathbf F_2\alpha_8\\
 D_3&\mathbf F_2\eta_5\\
 N&\displaystyle\left(\bigoplus_{q\text{ odd prime}}\mathbf F_2\psi_q\right)
          \oplus\mathbf F_2\eta_3.
\end{array}
\]
\end{corollary}

\begin{proof}
The first two rows follow from the exact scalar and joint quotient
maps. A character on $D_5$ killed on $D_9$ restricts on $D_8$ to
zero or $\eta_9$; the latter is impossible by the preceding
proposition. Its restriction therefore already kills $D_8$,
whose kernel on $D_5$ is the displayed span of $\eta_7,\alpha_8$.
No nonzero combination of those two characters extends to $D_3$,
so the $D_3$ kernel is unchanged. Similarly any invariant character
on $N$ killed on $D_9$ must kill $D_8$, since otherwise its
restriction to $D_5$ would extend $\eta_9$. The preceding kernel
formula for $N$ then gives the last row.
\end{proof}

